\chapter{CNN Architectures: VGG, ResNet, EfficientNet}
\label{ch:arch}

\chapmeta{Intermediate}{$\sim$4\,h (2\,h + 2\,h)}{EuroSAT}

\begin{objbox}
\begin{itemize}[leftmargin=1.2em,nosep]
  \item Trace the evolution from LeNet/AlexNet/VGG to ResNet and EfficientNet.
  \item Explain the degradation problem and how residual connections solve it.
  \item Implement a ResNet basic block and reason about gradient flow.
  \item Understand EfficientNet's compound scaling and the accuracy/efficiency frontier.
  \item Swap backbones in a classifier and compare parameters, speed and accuracy.
\end{itemize}
\end{objbox}

\section{A short architectural history}

LeNet-5 (1998) established the \texttt{conv--pool--FC} template. AlexNet (2012)
scaled it up with ReLU, dropout and GPU training, winning ImageNet by a wide
margin. VGG (2014) showed that \emph{depth} built from uniform $3\times3$
convolutions works well --- recall (Chapter~\ref{ch:conv}) that two stacked
$3\times3$ layers match a $5\times5$ receptive field with fewer parameters and
an extra non-linearity. But naively stacking ever more layers \emph{degraded}
accuracy, even on the training set --- not overfitting, but an optimisation
failure. ResNet (2015) fixed this and remains the default EO encoder.

\section{The degradation problem and residual connections}

A residual block learns a \emph{residual} $\mathcal{F}(\mathbf{x})$ added to its
input via a skip (identity) connection:
\begin{equation}
\mathbf{y} = \mathcal{F}(\mathbf{x}; \{W_i\}) + \mathbf{x}.
\label{eq:resid}
\end{equation}
If the optimal mapping for a block is close to the identity, it is far easier to
drive $\mathcal{F}\to 0$ than to make a stack of nonlinear layers represent the
identity exactly. Crucially, the skip provides a gradient highway: by the chain
rule the gradient reaching $\mathbf{x}$ contains an additive ``$+1$'' term,
\begin{equation}
\frac{\partial \mathcal{L}}{\partial \mathbf{x}}
= \frac{\partial \mathcal{L}}{\partial \mathbf{y}}
\left(1 + \frac{\partial \mathcal{F}}{\partial \mathbf{x}}\right),
\end{equation}
so gradients cannot vanish merely from depth. This lets networks train at 50,
101, even 152 layers. Figure~\ref{fig:resblock} shows the basic block.

\begin{figure}[h]
\centering
\fitwidth{
\begin{tikzpicture}[font=\small,>=Latex,node distance=8mm]
  \tikzstyle{op}=[draw,rounded corners,minimum width=26mm,minimum height=7mm,fill=eoblue!8]
  \node (x) {$\mathbf{x}$};
  \node[op,right=8mm of x] (c1) {Conv $3\times3$};
  \node[op,right=6mm of c1] (b1) {BN};
  \node[op,right=6mm of b1] (r1) {ReLU};
  \node[op,right=6mm of r1] (c2) {Conv $3\times3$};
  \node[op,right=6mm of c2] (b2) {BN};
  \node[draw,circle,right=10mm of b2,inner sep=1.5pt] (sum) {$+$};
  \node[op,right=8mm of sum,fill=eogreen!12] (r2) {ReLU};
  \node[right=8mm of r2] (y) {$\mathbf{y}$};
  \draw[->] (x)--(c1); \draw[->] (c1)--(b1); \draw[->] (b1)--(r1);
  \draw[->] (r1)--(c2); \draw[->] (c2)--(b2); \draw[->] (b2)--(sum);
  \draw[->] (sum)--(r2); \draw[->] (r2)--(y);
  \draw[->,eored,thick] (x.south) -- ++(0,-9mm) -| (sum.south)
        node[pos=0.25,below,black]{identity skip $\mathbf{x}$};
\end{tikzpicture}}
\caption{ResNet basic block (Eq.~\ref{eq:resid}). The identity skip (red) lets
the block default to a near-identity map and gives gradients an unobstructed
path, enabling very deep training. A $1\times1$ ``projection'' conv replaces the
skip when channel count or resolution changes.}
\label{fig:resblock}
\end{figure}

Deeper ResNets replace the two-$3\times3$ block with a \emph{bottleneck}
($1\times1\to3\times3\to1\times1$) that reduces and then restores channel
dimension, cutting cost while keeping capacity --- the building block of
ResNet-50/101.

\section{EfficientNet and compound scaling}

Networks can grow along three axes: depth (layers), width (channels) and input
resolution. EfficientNet (2019) showed these should be scaled \emph{together} in
fixed ratios governed by a single compound coefficient $\phi$:
\begin{equation}
\text{depth} \propto \alpha^{\phi},\quad
\text{width} \propto \beta^{\phi},\quad
\text{resolution} \propto \gamma^{\phi},\qquad
\alpha\beta^2\gamma^2 \approx 2,
\end{equation}
where increasing $\phi$ by one roughly doubles compute. Built on
mobile inverted-bottleneck (MBConv) blocks with squeeze-and-excitation,
EfficientNet-B0$\dots$B7 trace an accuracy/efficiency frontier that beats
ResNets at equal FLOPs --- attractive when EO inference must run cheaply at
scale.

\section{Swapping backbones}

In code, changing architecture is a one-line change because \texttt{torchvision}
and \texttt{timm} expose a uniform API:
\begin{lstlisting}
import torchvision.models as tv
m18 = tv.resnet18(weights=None, num_classes=10)
m50 = tv.resnet50(weights=None, num_classes=10)
import timm
eff = timm.create_model("efficientnet_b0", num_classes=10, in_chans=3)
\end{lstlisting}
Benchmarking on EuroSAT typically shows ResNet18 ($\sim$11\,M params) training
fastest, ResNet50 ($\sim$25\,M) slightly more accurate, and EfficientNet-B0
($\sim$5\,M) matching ResNet50 accuracy at a fraction of the parameters --- the
point of compound scaling. The same encoders reappear inside segmentation
decoders in Chapters~\ref{ch:seg}--\ref{ch:advseg}, which is exactly why we study
them now.

\begin{keybox}
\begin{itemize}[leftmargin=1.2em,nosep]
  \item VGG: depth from stacked $3\times3$; expensive FC head.
  \item Degradation $\ne$ overfitting; residual skips (Eq.~\ref{eq:resid}) solve it and ease gradient flow.
  \item Bottleneck blocks scale ResNet to 50/101/152 layers efficiently.
  \item EfficientNet scales depth/width/resolution jointly; best accuracy per FLOP.
  \item The classification encoder is reused as the segmentation encoder later.
\end{itemize}
\end{keybox}

\begin{practbox}
Notebook: \nbpath{chapter\_06\_cnn\_architectures/06\_practical\_cnn\_architectures.ipynb}.
\begin{enumerate}[leftmargin=1.4em,nosep]
  \item Implement a \texttt{BasicBlock} from scratch with the identity skip.
  \item Visualise gradient magnitude per layer \emph{with} and \emph{without} the
  skip to see the gradient highway.
  \item Tabulate parameters, forward-pass latency and accuracy for ResNet18/50 and
  EfficientNet-B0.
  \item Train ResNet18 vs EfficientNet-B0 on EuroSAT for a few epochs and compare.
\end{enumerate}
\end{practbox}

\begin{exbox}
\textbf{6.1 Bottleneck block.} Implement the $1\times1\to3\times3\to1\times1$
bottleneck and verify the parameter saving versus two $3\times3$ layers at equal
channels.\\[2pt]
\textbf{6.2 Depth ablation.} Build plain (no-skip) vs residual nets of increasing
depth; show that only the residual ones keep improving.\\[2pt]
\textbf{6.3 Width vs depth.} At a fixed parameter budget, compare a wider-shallow
and a narrower-deep network on EuroSAT.\\[2pt]
\textbf{Mini-project.} Write \texttt{build\_classifier(backbone, n\_classes)}
supporting ResNet18/50 and EfficientNet-B0, and a benchmark harness that outputs
a params-vs-accuracy plot.
\end{exbox}
